\documentclass[12pt]{amsart}
\usepackage[margin=1in]{geometry}
\pagestyle{plain}

\usepackage{amsmath,amssymb}
\usepackage{enumerate}
\usepackage[shortlabels]{enumitem}
\usepackage{graphicx}
\usepackage{xcolor}
\usepackage[T1]{fontenc}
\usepackage[parfill]{parskip}
\renewcommand{\baselinestretch}{1.1}

% Custom commands for clarity
\newcommand{\correct}[1]{{\color{blue}\textbf{CORRECT: #1}}}
\newcommand{\partcred}[1]{{\color{red}\textit{Partial Credit: #1}}}
\newcommand{\hint}{\textbf{Key Insight: }}
\newcommand{\trap}{\textbf{Common Trap: }}

\title{ESE 2030 QUIZZAM 2 : Solution Guide\\Fall 2025}
\author{Solution Key with Explanations}

\begin{document}
\maketitle

\section*{Overview}
This solution guide provides detailed explanations for each problem on Quiz 2, covering Weeks 3-5 material: quotient spaces, FTLA, change of basis, inner products, orthogonality, and QR decomposition.

\hrulefill

\section*{Problem 1: Coimage Dimension}
\textbf{Topic:} Coimage, quotient spaces, rank-nullity theorem

\correct{(A) 3}

\textbf{Explanation:} 
The coimage is defined as $\text{coim}(T) = V/\ker(T)$.

Using the Rank-Nullity Theorem:
\begin{itemize}
\item $\dim(V) = \text{rank}(T) + \dim(\ker(T))$
\item $9 = 3 + \dim(\ker(T))$
\item Therefore: $\dim(\ker(T)) = 6$
\end{itemize}

The dimension of the coimage:
$$\dim(\text{coim}(T)) = \dim(V) - \dim(\ker(T)) = 9 - 6 = 3$$

\hint The coimage "mods out" the kernel, so its dimension equals the rank.

\hrulefill

\section*{Problem 2: Orthonormal Set Properties}
\textbf{Topic:} Orthonormal sets, inner products

\correct{(A) Any linear combination has length $\sqrt{c_1^2 + c_2^2 + c_3^2}$}

\partcred{(D) True if $V$ is 3-dimensional and we can form a matrix}

\textbf{Explanation:} 
For an orthonormal set $\mathcal{S} = \{\mathbf{v}_1, \mathbf{v}_2, \mathbf{v}_3\}$:

$$\left\|c_1\mathbf{v}_1 + c_2\mathbf{v}_2 + c_3\mathbf{v}_3\right\|^2 = \sum_{i,j} c_ic_j\langle\mathbf{v}_i, \mathbf{v}_j\rangle$$

Since the vectors are orthonormal: $\langle\mathbf{v}_i, \mathbf{v}_j\rangle = \delta_{ij}$

Therefore: $\|c_1\mathbf{v}_1 + c_2\mathbf{v}_2 + c_3\mathbf{v}_3\|^2 = c_1^2 + c_2^2 + c_3^2$

Why other options are wrong:
\begin{itemize}
\item (B): They span a 3D subspace but $V$ could be larger
\item (C): $\|\mathbf{v}_1 + \mathbf{v}_2 + \mathbf{v}_3\|^2 = 1 + 1 + 1 = 3 \neq 1$
\item (D): Only true if we can form a matrix (requires $V = \mathbb{R}^3$)
\end{itemize}

\hrulefill

\section*{Problem 3: Operations Without Inner Product}
\textbf{Topic:} Inner product requirements vs. algebraic operations

\correct{(C) I and III only}

\partcred{(A) Recognizes basis checking doesn't need inner product}

\textbf{Explanation:} 
\begin{itemize}
\item \textbf{Statement I (TRUE):} Checking if vectors form a basis only requires:
  \begin{itemize}
  \item Linear independence (solve $c_1(x) + c_2(1+x^2) + c_3(1+x+x^2) = 0$)
  \item Spanning (verify we get all of $\mathcal{P}_2$)
  \item These are purely algebraic operations
  \end{itemize}
\item \textbf{Statement II (FALSE):} Orthogonality requires computing $\langle x, x^2 \rangle$, which needs an inner product definition
\item \textbf{Statement III (TRUE):} Finding dimension only requires checking linear independence via row reduction
\end{itemize}

\hint Geometric concepts (angles, orthogonality) need inner products; algebraic concepts (independence, dimension) don't.

\hrulefill

\section*{Problem 4: Change of Basis Coordinates}
\textbf{Topic:} Change of basis, coordinate vectors

\correct{(B) $\begin{pmatrix} 1 \\ -1 \\ 3 \end{pmatrix}$}

\textbf{Explanation:} 
Given: $\mathbf{v} = 2\mathbf{b}_1 - \mathbf{b}_2 + 3\mathbf{b}_3$

New basis: $\mathcal{C} = \{2\mathbf{b}_1, \mathbf{b}_2, \mathbf{b}_3\}$

We need to express $\mathbf{v}$ in terms of $\mathcal{C}$:
$$\mathbf{v} = 2\mathbf{b}_1 - \mathbf{b}_2 + 3\mathbf{b}_3 = 1 \cdot (2\mathbf{b}_1) + (-1) \cdot \mathbf{b}_2 + 3 \cdot \mathbf{b}_3$$

Therefore: $[\mathbf{v}]_{\mathcal{C}} = \begin{pmatrix} 1 \\ -1 \\ 3 \end{pmatrix}$

\trap The first coordinate changes because we're using $2\mathbf{b}_1$ as our new first basis vector.

\hrulefill

\section*{Problem 5: Similarity Invariants}
\textbf{Topic:} Similar matrices, invariant properties

\correct{(D) The column space}

\partcred{(E) Recognizes that most properties are preserved}

\textbf{Explanation:} 
Similar matrices $B = P^{-1}AP$ preserve:
\begin{itemize}
\item Rank (dimension of image)
\item Nullity (dimension of kernel)
\item Determinant: $\det(B) = \det(P^{-1})\det(A)\det(P) = \det(A)$
\item Trace, eigenvalues, characteristic polynomial
\end{itemize}

NOT preserved:
\begin{itemize}
\item Column space: Changes with basis transformation
\item Row space: Also changes
\item Specific eigenvectors (though eigenspaces are preserved)
\end{itemize}

Example: $I$ and $P^{-1}IP = I$ are similar but may have different column spaces if viewed as subspaces of $\mathbb{R}^n$ with $P \neq I$.

\hrulefill

\section*{Problem 6: Matrix Representations Under Basis Change}
\textbf{Topic:} Change of basis for linear transformations

\correct{(B) They have the same rank}

\textbf{Explanation:} 
$[T]_{\mathcal{B}}$ and $[T]_{\mathcal{C}}$ are similar matrices related by:
$$[T]_{\mathcal{C}} = P^{-1}[T]_{\mathcal{B}}P$$

where $P$ is the change of basis matrix. Similar matrices preserve:
\begin{itemize}
\item Rank (both equal $\dim(\text{im}(T))$)
\item Determinant
\item Trace
\item Eigenvalues
\end{itemize}

They do NOT have:
\begin{itemize}
\item Same column vectors (change with basis)
\item Same diagonal entries (except trace is preserved)
\item They represent the SAME transformation (not different ones)
\end{itemize}

\hrulefill

\section*{Problem 7: Cokernel Properties}
\textbf{Topic:} Cokernel, quotient spaces

\correct{(C) $\dim(\text{coker}(T)) = 2$ but it need not be isomorphic to $\ker(T)$}

\partcred{(A) Correct dimension but incomplete understanding}

\textbf{Explanation:} 
The cokernel is $\text{coker}(T) = W/\text{im}(T) = \mathbb{R}^5/\text{im}(T)$.

Dimension calculation:
$$\dim(\text{coker}(T)) = \dim(W) - \dim(\text{im}(T)) = 5 - 3 = 2$$

While both $\ker(T)$ and $\text{coker}(T)$ have dimension 2:
\begin{itemize}
\item $\ker(T)$ is a subspace of the domain $\mathbb{R}^5$
\item $\text{coker}(T)$ is a quotient of the codomain $\mathbb{R}^5$
\item There's no natural isomorphism between them
\item They're abstractly isomorphic as 2D vector spaces but not canonically
\end{itemize}

\hrulefill

\section*{Problem 8: QR Decomposition and Similarity}
\textbf{Topic:} QR decomposition, matrix multiplication

\correct{(E) $RQ$}

\textbf{Explanation:} 
Given $A = QR$ where $Q$ is orthogonal (so $Q^{-1} = Q^T$):

$$B = Q^{-1}AQ = Q^T(QR)Q$$

Simplifying:
$$B = (Q^TQ)RQ = IRQ = RQ$$

Note: This is the key step in the QR algorithm for eigenvalue computation!hich we will learn later!

\hrulefill

\section*{Problem 9: Norm in Orthonormal Basis}
\textbf{Topic:} Orthonormal bases, norm computation

\correct{(A) $30$}

\textbf{Explanation:} 
For $\mathbf{w} = 3\mathbf{v}_1 - 2\mathbf{v}_2 + \mathbf{v}_3 - 4\mathbf{v}_4$ with orthonormal basis $B$:

$$\|\mathbf{w}\|^2 = \langle\mathbf{w}, \mathbf{w}\rangle$$

Using orthonormality ($\langle\mathbf{v}_i, \mathbf{v}_j\rangle = \delta_{ij}$):
$$\|\mathbf{w}\|^2 = 3^2\|\mathbf{v}_1\|^2 + (-2)^2\|\mathbf{v}_2\|^2 + 1^2\|\mathbf{v}_3\|^2 + (-4)^2\|\mathbf{v}_4\|^2$$
$$= 9 + 4 + 1 + 16 = 30$$

\hint With orthonormal bases, the norm squared equals the sum of squared coefficients.

\hrulefill

\section*{Problem 10: QR Decomposition Properties}
\textbf{Topic:} Properties of Q and R matrices

\correct{(C) The columns of $Q$ form an orthonormal basis}

\textbf{Explanation:} 
In QR decomposition:
\begin{itemize}
\item $Q$ has orthonormal columns by definition  
\item $R$ is upper triangular
\item Diagonal entries of $R$ can be negative (depends on algorithm)
\item $R$ is invertible only if $A$ has full rank
\item For orthogonal $Q$: $Q^TQ^{-1} = Q^T(Q^T) = (Q^T)^2 \neq I$ in general
\end{itemize}

Why option (A) is false: Since $Q^{-1} = Q^T$ for orthogonal matrices, we get $Q^TQ^{-1} = (Q^T)^2$, which is NOT the identity matrix unless $Q^T$ is an involution (rare special case).

\hrulefill

\section*{Problem 11: Orthogonal Matrix Properties}
\textbf{Topic:} Properties of orthogonal matrices

\correct{(E) None of the above statements is false}

\textbf{Explanation:} 
For an orthogonal matrix $Q$ ($Q^TQ = I$):
\begin{itemize}
\item \textbf{(A) TRUE:} $\langle Qx, Qy \rangle = (Qx)^T(Qy) = x^TQ^TQy = x^Ty = \langle x, y \rangle$
\item \textbf{(B) TRUE:} Columns are orthonormal by definition
\item \textbf{(C) TRUE:} $(Q^T)^TQ^T = QQ^T = I$, so $Q^T$ is orthogonal
\item \textbf{(D) TRUE:} Each column has unit norm, so $|q_{ij}| \leq 1$ for all entries
\end{itemize}

All statements are true!

\hrulefill

\section*{Problem 12: Weighted Inner Product}
\textbf{Topic:} Weighted inner products, induced norms

\correct{(D) $\|f\| = \sqrt{\int_0^1 f(x)^2 e^x dx}$}

\textbf{Explanation:} 
Given the weighted inner product $\langle f, g \rangle = \int_0^1 e^x f(x)g(x) dx$:

The induced norm is:
$$\|f\| = \sqrt{\langle f, f \rangle} = \sqrt{\int_0^1 e^x f(x)f(x) dx} = \sqrt{\int_0^1 f(x)^2 e^x dx}$$

Why others are wrong:
\begin{itemize}
\item (A): This IS a valid inner product (weight $e^x > 0$)
\item (B): $\|1\|^2 = \int_0^1 e^x dx = e - 1$, so $\|1\| = \sqrt{e-1} \neq e$
\item (C): Orthogonality depends on the specific inner product
\end{itemize}

\hrulefill

\section*{Problem 13: Angle Between Linear Combinations}
\textbf{Topic:} Inner products, angle computation

\correct{(D) $\frac{2\|\mathbf{u}\|^2 - 3\langle\mathbf{u},\mathbf{v}\rangle}{\|\mathbf{u}\| \cdot \|\mathbf{w}\|}$}

\partcred{(C) Attempts correct approach but incorrectly simplifies}

\textbf{Explanation:} 
The cosine of the angle between $\mathbf{u}$ and $\mathbf{w} = 2\mathbf{u} - 3\mathbf{v}$ is:

$$\cos\phi = \frac{\langle\mathbf{u}, \mathbf{w}\rangle}{\|\mathbf{u}\|\|\mathbf{w}\|} = \frac{\langle\mathbf{u}, 2\mathbf{u} - 3\mathbf{v}\rangle}{\|\mathbf{u}\|\|\mathbf{w}\|}$$

Computing the numerator:
$$\langle\mathbf{u}, 2\mathbf{u} - 3\mathbf{v}\rangle = 2\langle\mathbf{u}, \mathbf{u}\rangle - 3\langle\mathbf{u}, \mathbf{v}\rangle = 2\|\mathbf{u}\|^2 - 3\langle\mathbf{u}, \mathbf{v}\rangle$$

\trap Don't confuse norms with inner products in the numerator!

\hrulefill

\section*{Problem 14: FTLA}
\textbf{Topic:} Quotient spaces, fundamental theorem of linear algebra

\correct{(B) They are isomorphic, and the isomorphism is induced by $T$ itself}

\partcred{(A) (small) Recognizes equal dimensions}

\textbf{Explanation:} 
The FTLA states:
$$\mathbb{R}^4/\ker(T) = \text{coim} \cong \text{im}(T)$$

Both spaces have dimension 2:
\begin{itemize}
\item Quotient: $\dim(\mathbb{R}^4/\ker(T)) = 4 - \dim(\ker(T)) = 4 - 2 = 2$
\item Image: $\dim(\text{im}(T)) = \text{rank}(T) = 2$
\end{itemize}

The isomorphism is given by: $[\mathbf{v}]_{\ker(T)} \mapsto T(\mathbf{v})$

This map is well-defined and bijective by the theorem.

\hrulefill

\section*{Problem 15: Angle-Preserving Transformations}
\textbf{Topic:} Orthogonal transformations, geometric properties

\correct{(B) $A^TA = I$ (orthogonal matrix)}

\textbf{Explanation:} 
For angle preservation, we need:
$$\cos\theta' = \frac{\langle T\mathbf{u}, T\mathbf{v}\rangle}{\|T\mathbf{u}\|\|T\mathbf{v}\|} = \frac{\langle \mathbf{u}, \mathbf{v}\rangle}{\|\mathbf{u}\|\|\mathbf{v}\|} = \cos\theta$$

This happens if and only if $T$ preserves inner products:
$$\langle T\mathbf{u}, T\mathbf{v}\rangle = \langle \mathbf{u}, \mathbf{v}\rangle$$

For matrix $A$: $\langle A\mathbf{x}, A\mathbf{y}\rangle = \mathbf{x}^TA^TA\mathbf{y} = \mathbf{x}^T\mathbf{y}$ requires $A^TA = I$.

Why others fail:
\begin{itemize}
\item (A): Symmetric matrices don't preserve angles in general
\item (C): Necessary but not sufficient
\item (D): Equal column lengths don't ensure orthonormality
\item (E): Skew-symmetric matrices don't preserve angles
\end{itemize}

\hrulefill

\section*{Problem 16: Scalar Multiplication Geometry}
\textbf{Topic:} Geometric properties of linear transformations

\correct{(C) Angles between vectors}

\textbf{Explanation:} 
For $T(\mathbf{x}) = 5\mathbf{x}$:

\begin{itemize}
\item Inner products: $\langle 5\mathbf{u}, 5\mathbf{v}\rangle = 25\langle\mathbf{u}, \mathbf{v}\rangle$ (scaled by 25)
\item Lengths: $\|5\mathbf{u}\| = 5\|\mathbf{u}\|$ (scaled by 5)
\item Distances: $\|5\mathbf{u} - 5\mathbf{v}\| = 5\|\mathbf{u} - \mathbf{v}\|$ (scaled by 5)
\item Angles: $\cos\theta' = \frac{25\langle\mathbf{u}, \mathbf{v}\rangle}{5\|\mathbf{u}\| \cdot 5\|\mathbf{v}\|} = \frac{\langle\mathbf{u}, \mathbf{v}\rangle}{\|\mathbf{u}\|\|\mathbf{v}\|} = \cos\theta$  
\end{itemize}

\hint Uniform scaling preserves angles but not distances or inner products.

\hrulefill

\section*{Problem 17: Coimage Equivalence Classes}
\textbf{Topic:} Equivalence classes in quotient spaces

\correct{(D) $T(\mathbf{u}) = T(\mathbf{v})$}

\partcred{(E) True but trivially so, not the defining condition}

\textbf{Explanation:} 
Two vectors $\mathbf{u}, \mathbf{v}$ are in the same equivalence class of $\text{coim}(T) = V/\ker(T)$ if and only if:

$$\mathbf{u} - \mathbf{v} \in \ker(T)$$

This is equivalent to:
$$T(\mathbf{u} - \mathbf{v}) = \mathbf{0}$$
$$T(\mathbf{u}) = T(\mathbf{v})$$

Option (E) states $T(\mathbf{u}) - T(\mathbf{v}) \in \text{im}(T)$, which is always true but doesn't characterize equivalence.

\hrulefill

\section*{Problem 18: Inner Product Requirements}
\textbf{Topic:} Operations requiring vs. not requiring inner products

\correct{(B) Computing the adjoint $T^*$}

\partcred{(E) Matrix representation is easier with inner product but doesn't require it}

\textbf{Explanation:} 
The adjoint $T^*$ is defined by the relation:
$$\langle T\mathbf{v}, \mathbf{w}\rangle = \langle\mathbf{v}, T^*\mathbf{w}\rangle$$

This definition explicitly requires an inner product.

Operations NOT requiring inner products:
\begin{itemize}
\item Invertibility: Check if $\ker(T) = \{\mathbf{0}\}$
\item Finding kernel: Solve $T\mathbf{v} = \mathbf{0}$
\item Coimage basis: Purely algebraic quotient construction
\item Matrix representation: Only needs choice of bases
\end{itemize}

\hrulefill

\section*{Problem 19: Identifying Orthogonal Matrices}
\textbf{Topic:} Orthogonal matrices, computational verification

\textbf{CRITICAL ERROR - PROBLEM IS BROKEN}

The problem asks "Which of the following matrices is orthogonal?" but THREE matrices are orthogonal!

\textbf{Actual results:}
\begin{itemize}
\item \textbf{$Q_1$ is ORTHOGONAL} 
\item \textbf{$Q_2$ is ORTHOGONAL} 
\item \textbf{$Q_3$ is NOT orthogonal} (columns 1 and 3 not orthogonal)
\item \textbf{$Q_4$ is ORTHOGONAL} 
\end{itemize}

Option (E) "None of the above is orthogonal" is FALSE.

\textbf{This problem has no unique correct answer.} Options A, B, and D are all correct. Apologies from prof-g.

For grading: Full credit to students who chose A, B, or D, and no credit for C or E (which is factually incorrect).

\hrulefill

\hrulefill

\section*{Problem 20: Document Vector Angles}
\textbf{Topic:} Vector addition and angles

\correct{(C) These angles are both strictly less than $30°$}

\partcred{(A) Might be exactly $15°$ in special cases}
\partcred{(E) Can determine bounds without knowing lengths}

\textbf{Explanation:} 
Given $\mathbf{d}_3 = \mathbf{d}_1 + \mathbf{d}_2$ with $\theta_{1,2} = 30°$:

By the parallelogram law, $\mathbf{d}_3$ lies in the interior of the angle between $\mathbf{d}_1$ and $\mathbf{d}_2$ (when viewed from origin).

Both angles $\theta_{1,3}$ and $\theta_{2,3}$ must be:
\begin{itemize}
\item Strictly less than $30°$ (the original angle)
\item Strictly greater than $0°$ (vectors point in same general direction)
\item Equal to $15°$ only if $\|\mathbf{d}_1\| = \|\mathbf{d}_2\|$
\end{itemize}

\hint Vector addition creates a "weighted average" direction between the two vectors.

\hrulefill

\section*{Problem 21: Matrix Representation of Differentiation}
\textbf{Topic:} Matrix representations of linear operators

\correct{(A) The columns are coordinate vectors of $T(1), T(x), T(x^2)$ with respect to $\mathcal{C}$}

\partcred{(E) Related idea but incorrect conclusion}

\textbf{Explanation:} 
Computing $[T]_{\mathcal{B}}^{\mathcal{C}}$:

\begin{itemize}
\item $T(1) = 0 = 0 \cdot 1 + 0 \cdot x$ → column 1: $\begin{pmatrix}0\\0\end{pmatrix}$
\item $T(x) = 1 = 1 \cdot 1 + 0 \cdot x$ → column 2: $\begin{pmatrix}1\\0\end{pmatrix}$
\item $T(x^2) = 2x = 0 \cdot 1 + 2 \cdot x$ → column 3: $\begin{pmatrix}0\\2\end{pmatrix}$
\end{itemize}

Therefore: $[T]_{\mathcal{B}}^{\mathcal{C}} = \begin{bmatrix} 0 & 1 & 0 \\ 0 & 0 & 2 \end{bmatrix}$

The matrix is $2 \times 3$ (not square), and changing basis won't change rank.

\hrulefill

\section*{Problem 22: Adjoint of Differentiation}
\textbf{Topic:} Adjoint operators, inner products

\correct{(E) For all $p \in \mathcal{P}_3$ and $q \in \mathcal{P}_2$: $\int_0^1 p'(x)q(x)dx = \int_0^1 p(x)(D^*(q))(x)dx$}

\partcred{(C) Related to integration by parts but vague}

\textbf{Explanation:} 
The defining property of the adjoint is:
$$\langle D(p), q \rangle = \langle p, D^*(q) \rangle$$

With our inner product $\langle f, g \rangle = \int_0^1 f(x)g(x)dx$:
$$\int_0^1 p'(x)q(x)dx = \int_0^1 p(x)(D^*(q))(x)dx$$

This is exactly what option (E) states.

Note: $D^*$ is NOT simply antidifferentiation due to boundary terms from integration by parts. The actual formula involves both antiderivatives and boundary corrections.

\hint The adjoint relationship is the fundamental definition—everything else follows from it.

\hrulefill

\section*{Key Takeaways for Quiz 2}

\begin{enumerate}
\item \textbf{Quotient Spaces:} Remember dim$(V/W) = \dim(V) - \dim(W)$
\item \textbf{Orthonormal Bases:} Dramatically simplify norm and inner product calculations
\item \textbf{QR Decomposition:} $Q$ has orthonormal columns; $R$ is upper triangular
\item \textbf{FTLA:} $V/\ker(T) \cong \text{im}(T)$
\item \textbf{Inner Products:} Required for geometry (angles, adjoints) but not algebra (rank, independence)
\item \textbf{Change of Basis:} Similar matrices preserve rank, determinant, eigenvalues—but not column spaces
\end{enumerate}

\end{document}